% critique.tex — f162v, vv. 1–42 : commentaire critique
% Articles tirés du commentaire antérieur de l'auteur (texte et notes conservés) ; répartition : repartition.md
\cc{16–8}{πρῶτα ἐμηχανήσατο ... ἀκοὴν δὲ χοάνης ὦτα διετετρήνατο.}{Pour l'explication de ce passage diversement retouché au gré des éditeurs, je ne peux que renvoyer au brillant exercice d'exégèse d'Alphonse Willems. (Ses notes sur les \italique{Thesmophories} sont reproduites en annexe, celle relative à ce passage en est la première. À propos d'Alphonse Willems, \italique{confer} l'Introduction.) Willems construit ainsi la phrase, en accord avec la sobre indication du scholiaste (λείπει ὡς) : διετετρήνατο δὲ ἀκοὴν (ὡς) χοάνης ὦτα.}
\cc{22–4}{πῶς ἂν οὖν ... προσμάθοιμι χωλὸς εἶναι τὼ σκέλει;}{(Voir là-dessus la \nieme{3} note d'Alphonse Willems, reproduite en annexe, qui justifie la suppression du vers 23 -- une glose explicative du suivant.)}
\cc{30–3}{ποῖος οὗτος Ἁγάθων; ... οὐχ ἑόρακας πώποτε.}{La répartition des répliques de ce passage fluctue d'une édition à l'autre. Tâchons d'y voir plus clair !\\
À l'instar du manuscrit de Ravenne, je laisse à Mnésiloque le début du vers 31 : ἔστιν τις Ἀγάθων; 
La clef de compréhension est la suivante : lisons :\citComm{ἔστιν τις ἀγαθῶν;}Mnésiloque joue sur le nom d'Agathon ; \latin{nomen omen}, demande-t-il : Agathon est-il bien l'\italique{honnête homme} que suggère son nom ? Or, pour une nature rustre comme celle de Mnésiloque, un bon-homme se doit d'être un solide gaillard, à la peau bronzée par le travail en extérieur et non pas au teint fromage-blanc comme ces désœuvrés de fémoïdes ! 
Il continue sur sa lancée : μῶν ὁ μέλας ὁ καρτερός; « n'est-ce pas le noiraud, le costaud ? »
(Agathon est en réalité tout le contraire d'un rustaud ; Aristophane donne ainsi facilement à rire à son public aux dépens de l'original bien connu.)\\
Le vers suivant est écrit ainsi dans le manuscrit de Ravenne :\citComm{οὔκ ἀλλ᾽ ἕτερός τις --- οὐχ \textbf{ἑόρακα} πώποτε}Le signe mis pour signifier un changement d'interlocuteur\nc{Non un tiret, \italique{paragraphos}, ni un double point ` : ', dont l'usage est courant dans notre manuscrit lorsque le changement d'interlocuteur intervient au milieu d'un vers, mais en l'occurrence une indication de personnage : μ\auDessus{ν}{Η}.} ayant par la suite été gratté (\latin{ut videtur}).
L'entièreté du vers est effectivement à restituer à Euripide : une scholie écrite juste au-dessus du vers l'atteste : \citComm{\textbf{ἐπειδὴ} οὐ πάλαι ἤρξατο διδάσκειν, ἀλλὰ τρισὶν πρὸ τούτων ἔτεσιν.\\
\textbf{Puisqu}'il n'a pas commencé récemment de faire jouer des pièces, mais depuis trois ans. [sept, en réalité, et selon la façon grecque de compter]}
Cette scholie est évidemment là pour expliquer l'étonnement d'Euripide : οὐχ ἑόρακα\textbf{ς} πώποτε; « Tu ne l'as donc jamais vu ?! [alors que cela fait trois années qu'il monte des pièces] » ; le scholiaste ne l'aurait pas ainsi rédigée (en commençant par ἐπειδὴ...) si, au vers 32, le `οὐχ ἑόρακα πώποτε' était prononcé par Mnésiloque, « je ne l'ai jamais vu ».\\
Cette interprétation, outre qu'elle est anciennement corroborée par la scholie, est aussi celle qui permet la lecture la plus drôle : οὐχ ἑόρακας πώποτε; --- μῶν ὁ δασυπώγων; --- οὐχ ἑόρακας πώποτε, « Tu ne l'as jamais vu ?! --- N'est-il pas barbu ? --- Tu ne l'as jamais vu !.. », bien meilleure par le rythme et la répétition déconvenue d'Euripide que l'échange premièrement transcrit sur le manuscrit de Ravenne où Mnésiloque commençait par avouer son ignorance : (Μνη.) οὐχ ἑόρακα πώποτε. μῶν ὁ δασυπώγων; --- (Εὐρ.) οὐχ ἑόρακας πώποτε.}
\cc{35}{καὶ μὴν βεβίνηκας σύ γ᾽, ἀλλ᾽ οὐκ οἶσθ᾽ ἴσως.}{Le ton est inhabituel -- voire plutôt choquant -- dans la bouche d'Euripide.
Qu'à cela ne tienne : puisqu'il semble d'autre part qu'Agathon tourne le dos lorsqu'il accorde ses faveurs\nc{Voir le vers 258 et sa scholie ; Agathon se livrerait au travestissement dans l'interlope \italique{Athens by night} pour y être éprouvé à la manière d'une femme -- mais le commentaire déluré du scholiaste n'est peut-être qu'une mauvaise interprétation du présent vers 35 adjointe à une forte dose d'imagination lubrique. Ce qui est sûr, c'est qu'au vers 205, Agathon dit craindre s'il était reconnu d'être lynché pour avoir chapardé nuitamment les œuvres de la Cypris : les Athéniennes voient en lui un concurrent irrégulier.}, nos commentateurs en concluent qu'il est logique que Mnésiloque ne le connaisse pas -- sauf éventuellement au sens biblique... de Sodome, s'entend. 
Et c'est le grave et docte Euripide qui colporterait ce genre de ragots. Y a-t-il besoin de commenter ce genre de commentaires ?.. je précise : ces commentaires \italique{modernes} ; car le βεβίνηκας du \italique{Ravennas} n'a pas laissé derrière lui la moindre scholie ancienne ; chose étonnante lorsqu'on sait que les anciens commentateurs étaient presqu'autant portés au scabreux que nos modernes. Il faut croire que le texte original d'Aristophane, l'\allemand{Urtext} des \italique{Thesmophories}, comportait une tout autre leçon : une banalité sur laquelle aucun commentateur ancien ne se sera arrêté.
Las !.. on ne détruit que ce que l'on remplace ; à défaut d'avoir une proposition alternative solidement étayée, il faut nous résoudre à \latin{crux}-désespérer...}
